\fontsize{12pt}{14pt} \selectfont
\begin{doublespace}
\chapter{Design}

The design of Houdini is guided by a set of core principles aimed at enabling reliable, large-scale testing of container security confinement. At its core, Houdini seeks to ensure \textbf{reproducibility} and \textbf{determinism}, so that each test yields the same result when run under identical conditions, and behaves consistently across time and environments. It enforces strong \textbf{host isolation} by executing potentially dangerous container workloads inside tightly controlled virtual machines, preventing any impact on the host system. The framework is built to support \textbf{flexibility across configurations}, allowing users to test a wide range of Docker flags, Linux capabilities, AppArmor profiles, and runtime versions. Despite its serialized test execution, Houdini is designed with \textbf{scalability} in mind through automation and declarative input modeling, enabling large test matrices to be evaluated with minimal manual effort. It also emphasizes \textbf{performance} by reducing overhead through lightweight virtual machines and precompiled binaries. Finally, Houdini prioritizes \textbf{transparency and debuggability} by logging all configuration inputs and test results in a structured format, enabling users to trace and interpret outcomes with confidence. The following sections describe how each of these design goals is achieved through Houdini's architecture and implementation.



\section{System Architecture}

Houdini is composed of several coordinated components that work together to execute container security tests in a controlled and reproducible manner. At the top level, the process is initiated by a host-side controller, a Python-based command-line interface (CLI)—which is responsible for parsing user input, defining the test matrix, and orchestrating the execution of each test case. Users interact with Houdini through this CLI to specify which trick to run, along with the desired combinations of container runtime versions, kernel versions, and configuration parameters such as capabilities, namespaces, and AppArmor profiles.

To facilitate communication between the host and the virtual test environment, Houdini runs a lightweight Flask server on the host. This server exposes a simple HTTP API that allows the host to send configuration data and trick logic into the virtual machine (VM) and retrieve results once execution is complete. The Flask server acts as a bridge between the Houdini CLI and the internal logic running inside the guest VM.

The actual test environment is provided by a QEMU-based virtual machine. This VM is constructed using a Buildroot-generated Linux image, which provides a minimal and deterministic guest operating system. Each test is launched in a fresh instance of this VM to ensure isolation, reproducibility, and protection for the host system. Even if a container escape occurs or a trick exploits kernel-level vulnerabilities, the damage is contained within the VM and cannot affect the host.

Inside the virtual machine, a component known as the \textit{trick listener} runs as a persistent background process. This service is responsible for receiving trick execution requests from the Flask server, parsing the incoming configuration, and launching the container under test with the specified parameters. The guest VM includes a version-pinned installation of Docker, which is used to start the container with the exact runtime environment defined by the test configuration. This includes setting flags such as \texttt{--pid=host}, applying AppArmor profiles, granting or dropping specific Linux capabilities, and mounting filesystems.

Once the container is running, the trick execution logic is carried out inside the container. This logic is defined by the user in a separate script—typically written in Python or Bash—that is copied into the container at build time. The script performs the security test (e.g., attempting to read a protected file, send a signal to a host process, or escalate privileges) and outputs a success or failure message depending on the observed behavior. After execution, the trick listener collects the results and sends them back to the host through the Flask server, where they are logged and stored for further analysis.

This architecture ensures strict separation between the host and guest environments, deterministic test behavior, and support for a wide variety of container configurations. By isolating execution within a lightweight VM and decoupling the configuration, orchestration, and evaluation layers, Houdini provides a robust and extensible platform for empirical container security testing.



\section{Trick Lifecycle and Execution Flow}

Each test in Houdini is structured around a unit called a \textit{trick}, which encapsulates a container security scenario that is to be evaluated. A trick consists of three core files: a configuration file, a Dockerfile, and an executable script. The configuration file defines how the container should be launched, including parameters such as Linux capabilities, namespace modes, AppArmor profiles, seccomp filters, volume mounts, and user privileges. This file also specifies the Dockerfile to use and the path to the trick implementation. The Dockerfile defines the build environment of the container, installs any required dependencies, and copies the trick script into the container image. Finally, the trick script itself—typically written in Bash or Python—contains the logic of the test. It attempts a specific action (such as reading \texttt{/etc/shadow}, sending a signal to a host process, or escalating privileges) and emits a clear “pass” or “fail” output depending on whether the action succeeds.

The trick lifecycle begins on the host, where the user selects a trick and provides a set of configurations to evaluate. The Houdini CLI sends this information to the Flask server, which forwards the request into the virtual machine via an HTTP API. Inside the VM, the trick listener receives the configuration and prepares the Docker container. Docker is used to build the container image using the specified Dockerfile, then run the container with the flags defined in the configuration file. This ensures that the environment within the container is exactly as described, with version-pinned binaries and well-defined privilege boundaries.

Once the container is launched, the trick script executes automatically inside the container. During execution, the script performs its predefined test logic and reports its outcome by printing a success or failure message to standard output. The trick listener monitors the container’s output and collects the final result, along with any standard error messages and the container's exit code. These results are then sent back to the Flask server on the host, which tags them with the relevant configuration metadata and stores them in a structured log format for later analysis.

To guarantee consistency and isolation, each trick is executed in a freshly booted QEMU virtual machine that uses a clean Buildroot-generated image. This eliminates state sharing between tests and prevents contamination from previous runs. No background services, daemons, or user-modifiable system components are present in the guest image, further ensuring that results are solely attributable to the test configuration. Because no persistent state is carried between executions, the test environment is fully reset each time, enabling strict reproducibility and simplifying debugging. This design guarantees that the observed behavior reflects the actual impact of the container configuration under test, free from side effects introduced by prior tests or host-level interactions.


\section{Version Pinning}

To ensure reproducibility and enable systematic regression testing, Houdini enforces strict control over the versions of all critical software components involved in test execution. Container behavior can change significantly between versions of Docker, runc, or the Linux kernel, even when the container configuration itself remains unchanged. To eliminate variability caused by underlying software drift, Houdini relies on a version pinning strategy that explicitly fixes the versions of all major runtime components.

Each test case specifies exact versions of Docker, runc, and the Linux kernel. These components are precompiled and stored in a local cache, ensuring that the same binaries are used across repeated runs. For example, if a trick is meant to evaluate a vulnerability present in runc 1.0.0-rc2 but fixed in 1.1.12, Houdini will use the exact corresponding runc binary for each test case without relying on package managers or internet downloads. Similarly, different versions of Docker (e.g., 20.10.3, 23.0.1, or 27.0.3) can be tested across the same trick logic to assess behavioral drift or identify regressions.

These precompiled binaries are injected into the guest VM using an overlay filesystem. Rather than rebuilding the entire root filesystem for each test, Houdini layers the required binaries over the static Buildroot image. This technique allows different test cases to share a common base environment while customizing their runtime stack independently. Overlay injection ensures that the guest VM sees only the specified versions of the relevant tools, regardless of what is installed on the host.

The guest operating system itself is built using Buildroot, a minimalist Linux build system that allows fine-grained control over the contents of the root filesystem. Houdini uses Buildroot to construct deterministic VM images that contain only the essential components needed for trick execution. This eliminates potential sources of nondeterminism that are common in general-purpose distributions, such as automatic updates, dependency resolution, and unpredictable background services.

By combining precompiled runtime components, overlay injection, and deterministic Buildroot-generated VMs, Houdini creates a fully controlled execution environment. This infrastructure ensures that any changes in trick behavior can be confidently attributed to intentional differences in configuration or runtime version—not to hidden changes in the surrounding software stack. It also enables robust regression testing, where the same trick is executed across a range of versions to observe how security guarantees evolve over time. This approach was particularly useful in the evaluation of CVE-2024-21626, where Houdini revealed that the vulnerability was present in runc versions earlier than those initially reported in the official CVE disclosure.


\section{Test Matrix Design}

Houdini is designed to support systematic and large-scale container security testing through a declarative test matrix model. Instead of requiring users to write custom scripts for each test case, Houdini allows users to define a matrix of configuration options that are automatically expanded into individual test cases. This approach simplifies the process of specifying and managing complex experiments, reduces manual effort, and improves coverage across different isolation and privilege configurations.

Each trick execution is defined by a set of configuration "knobs" that control how the container is launched. These knobs include Linux capabilities (e.g., \texttt{CAP\_SYS\_ADMIN}, \texttt{CAP\_SYS\_PTRACE}), namespace modes (such as \texttt{--pid=host} or \texttt{--net=host}), AppArmor profiles (e.g., \texttt{unconfined} or \texttt{docker-default}), seccomp profiles, user identity (root vs. non-root), and Docker flags like \texttt{--privileged}, \texttt{--read-only}, or volume mounts. Users specify the set of values they wish to test for each knob using a simple configuration format, typically YAML or JSON.

Once specified, Houdini automatically generates the Cartesian product of all parameter combinations, pairing each with the target trick and runtime version. For example, if a user wants to test a single trick across 3 Docker versions, 2 AppArmor profiles, and 2 sets of capabilities, Houdini will automatically generate 12 distinct test cases (3 × 2 × 2), each representing a unique configuration. This process eliminates the need for manual duplication and ensures that all combinations are systematically evaluated.

This declarative model significantly enhances scalability. Rather than managing each test manually, users can define broad sets of scenarios with a few lines of configuration. Houdini’s internal scheduler then handles the execution of each test case in a clean, isolated virtual machine, capturing results independently. This approach enables researchers to explore how different isolation settings interact, detect edge cases, and identify configurations that break expected security boundaries.

By abstracting test case specification into a declarative matrix, Houdini provides a scalable and maintainable framework for container security testing. It empowers users to conduct extensive coverage experiments with minimal manual configuration, while preserving precision and reproducibility in each individual test.


\section{Isolation and Safety Mechanisms}

Houdini is explicitly designed to safely execute potentially dangerous or privilege-escalating container workloads without risking the security or stability of the host system. To achieve this, it enforces strict isolation through a combination of virtual machine separation, containment policies, and runtime safeguards.

All trick executions occur inside a QEMU-based virtual machine, which serves as a hardware-level boundary between the test environment and the host. Unlike approaches that rely on Linux containers alone, this VM-level separation guarantees that even if a container successfully breaks out of its sandbox, it remains confined within the virtualized guest. This defense-in-depth model protects the host from both known and unknown container escape vulnerabilities, ensuring that test payloads cannot influence or damage the system running Houdini itself.

Houdini further avoids any direct linkage between the host and the guest environment. No host directories are mounted into the guest, and no namespaces are shared between the host and VM. Containers are launched entirely within the guest’s user, mount, PID, and network namespaces. This means that containers inside the VM have no visibility into host processes, filesystems, or resources. Features such as \texttt{--pid=host} only apply within the context of the VM, and never bridge to the real host.

To empirically verify that host access is not possible, Houdini incorporates a canary mechanism. During VM boot, a simulated host-accessible path is defined (e.g., a file or directory that mimics a mounted host volume). The trick under test may attempt to write to this path or read from it as part of the test logic. If the container can successfully access this canary, it indicates a misconfiguration or a potential escape vector. In practice, Houdini’s isolation model prevents such access unless the VM itself has been misconfigured, providing an additional layer of runtime validation.

Together, these mechanisms enforce strong host isolation and make Houdini a safe platform for evaluating security-critical container behavior. By containing tests within a tightly scoped virtual environment, it ensures that no trick—malicious or accidental—can compromise the integrity of the host system.


\section{Motivation for Virtualized Execution}

A central architectural decision in Houdini is to execute all container tests inside a lightweight virtual machine, rather than directly on the host system. This approach is primarily motivated by the need for environmental control, cross-system reproducibility, and kernel-level flexibility.

First, running inside a VM ensures a **clean and deterministic environment**. Houdini uses a minimal Buildroot-based Linux system as its virtualized testbed, which eliminates the variability and complexity commonly found in general-purpose operating systems. This avoids interference from background services, user-installed software, or automatic updates, thereby reducing noise in test outcomes and ensuring that any observed behavior is solely due to the container configuration under test.

Second, virtualization enables **reproducibility across different machines and over time**. Because the virtual machine is booted from a fixed kernel image and root filesystem, the entire runtime environment remains identical regardless of the underlying host. This makes it possible to compare results across test runs and systems with confidence, knowing that the execution context has not drifted. Such consistency is essential for empirically validating container security guarantees, especially when fine-grained differences in runtime behavior may influence security outcomes.

Third, the virtual machine setup grants Houdini the ability to **precisely control the kernel version** used during testing. This is a key capability for evaluating the evolution of confinement mechanisms across different Linux releases. Houdini can boot any custom or historical kernel version without modifying the host system, enabling security researchers to assess whether a given trick succeeds or fails as a function of kernel changes alone. This capability is essential for regression testing and understanding the security implications of kernel-level patches and configuration options.

Together, these properties make VM-based execution a foundational component of Houdini’s design, enabling systematic, consistent, and historically-aware container security testing.


\section{Automation and Execution Pipeline}

Houdini is designed to operate as a fully automated and repeatable test execution framework. Once a user specifies the desired trick and configuration matrix, the system takes over and handles all stages of test orchestration without requiring further manual intervention. This fully unattended execution pipeline reduces the user burden and enables large-scale evaluations to proceed with minimal effort.

Each test case is executed in isolation within its own QEMU virtual machine. The VMs are launched one at a time in a serialized fashion—Houdini does not currently support parallel execution. While this limits overall throughput, it offers a crucial benefit: predictability. Serial execution ensures that each test runs in a clean environment free from interference by other tasks, removing concurrency as a source of nondeterminism.

To enforce repeatability, Houdini initializes a fresh VM for every single test case. The guest image is regenerated using a minimal Buildroot system, and no state from previous runs is carried over. This guarantees that the environment is identical each time a trick is executed. Whether running a test for the first time or rerunning it weeks later with the same configuration, the outcome will be determined solely by the logic of the trick and the versions of the software under test—not by stale filesystem state, cached data, or unpredictable side effects.

This strict separation and automation pipeline allow users to queue hundreds or thousands of test cases and walk away, confident that results will be recorded consistently and reliably. Logs are saved automatically, configuration metadata is attached to every result, and errors are surfaced clearly for debugging. In effect, Houdini abstracts away the operational complexity of container security testing, allowing researchers and engineers to focus on analysis rather than setup and execution.


\section{Transparency and Debuggability}

Houdini is designed not only to execute tests reliably, but also to make results easy to interpret and debug. To support transparency and effective diagnosis, the framework implements structured logging, rich metadata tagging, and deterministic evaluation criteria for every test.

Each test execution generates a structured log file that includes the container's standard output, standard error, exit code, and a timestamped record of the configuration parameters used. These logs are stored in a standardized format and organized by test ID, making it straightforward to trace the provenance of any result. Whether a trick passed, failed, or encountered an unexpected runtime error, the corresponding log provides clear visibility into what occurred during execution.

Houdini also tags every test result with its full configuration metadata. This includes the container runtime version, kernel version, Docker flags, capabilities, namespace modes, and any AppArmor or seccomp profiles in effect. These tags are recorded alongside the test result, enabling automated analysis, filtering, and regression tracking. For example, a user can easily compare how the same trick behaves under different runc versions or with different namespace settings, based on the metadata attached to each result.

While the default mode is fully automated, Houdini also supports optional interactive inspection. Users can choose to pause execution at specific stages, such as before launching the container or after the trick completes. This allows them to manually inspect the state of the VM or container, check file system contents, or rerun the trick interactively for more detailed debugging.

Finally, each trick is written to emit a deterministic pass or fail outcome based on runtime observations. For instance, a trick might check for the presence of a specific file, the ability to signal another process, or the success of a privilege escalation attempt. These conditions are embedded in the trick logic itself, ensuring that the result of each test is consistent and reproducible across runs. By combining clear pass/fail semantics with detailed logs and metadata, Houdini provides a transparent and developer-friendly environment for security evaluation.

\section{Security Design}

A central aspect of Houdini's architecture is its emphasis on safe execution of potentially harmful container workloads. The primary security boundary is enforced at the level of a QEMU virtual machine, which isolates each test from the host system using hardware-level virtualization. This boundary is sufficient for Houdini’s threat model, which assumes that container workloads may attempt to exploit kernel vulnerabilities or escalate privileges within the guest. By design, no test has direct access to the host’s filesystem, namespaces, or process tree. Even if a trick succeeds in breaking out of its container, it only reaches the VM environment, which is ephemeral and isolated. Furthermore, no host-mounted volumes or devices are exposed to the guest, and Houdini does not rely on nested virtualization or host-provided namespaces. This model ensures that trick execution remains safely contained, protecting both the host and other tests from compromise.

\section{Extensibility}

Houdini is designed to be modular and extensible, allowing security researchers to easily add new tricks, runtimes, or configuration parameters. Each trick is defined by a set of modular files: a Dockerfile, a configuration file, and a script. Adding a new trick simply requires placing these files in a new subdirectory and registering the trick in the Houdini CLI. Similarly, new runtime versions—such as a different version of Docker, runc, or the Linux kernel—can be added to Houdini’s local cache, after which they become available as selectable options in the test matrix. While the current implementation focuses on Docker and runc, the system is structured in a way that would allow integration of other OCI-compliant runtimes like \texttt{crun} or \texttt{youki} with relatively few changes. This modularity ensures that Houdini can evolve alongside the container ecosystem and be used to evaluate future attack surfaces or isolation strategies with minimal developer effort.

\section{Buildroot Customization and Guest Environment Design}

Houdini's guest virtual machine is built using Buildroot, a minimalist embedded Linux build system that allows tight control over the software included in the image. This customization is critical for maintaining determinism, minimizing attack surface, and reducing boot time. The Buildroot image includes only the essential components required for container execution and trick evaluation. These include a statically linked BusyBox shell, the Docker daemon and CLI (injected via overlay), a Python interpreter, curl for HTTP communication, and a lightweight init system. 

Deliberate trade-offs are made to exclude non-essential services such as SSH, udev, systemd, or networking daemons, which could introduce nondeterministic behavior or additional complexity. Instead, the VM is initialized with a fixed set of devices and predictable kernel logs, enabling stable trick execution. Statically linking utilities where possible simplifies debugging by avoiding dynamic library mismatches and dependency issues across test runs. The result is a clean, lightweight VM that boots quickly and behaves identically across repeated invocations, serving as a robust platform for empirical security testing.


\section{Configuration Validation and Sanity Checks}

Before executing any test, Houdini performs a series of pre-run validation checks to ensure the integrity and completeness of the user-provided configuration. These include verifying that all required trick files are present (configuration file, Dockerfile, and script), that the specified Docker and runc versions exist in the local binary cache, and that configuration parameters such as capability names, AppArmor profiles, and namespace flags are syntactically valid.

The CLI also checks for duplicate or conflicting configurations, such as attempting to run a test with both \texttt{--privileged} and a restricted seccomp profile. If such conflicts are detected, the system emits a warning and skips the invalid case. Users are also warned if required components (e.g., QEMU, KVM, or Buildroot toolchain) are not installed or misconfigured on the host system.

These validation steps are critical for preventing wasted test runs, simplifying troubleshooting, and maintaining the integrity of large-scale evaluations. By catching errors early in the pipeline, Houdini ensures that only well-formed and meaningful test cases are executed.

\section{Why QEMU Virtual Machines Are Used}
To ensure safe, reproducible, and flexible testing of container security configurations, Houdini executes all experiments within QEMU-based virtual machines (VMs). This design choice is motivated by several technical and practical advantages offered by QEMU over alternative virtualization platforms like VirtualBox or VMware.

First and foremost, QEMU supports both full system emulation and hardware-assisted virtualization via KVM (Kernel-based Virtual Machine), allowing it to run unmodified Linux kernels and minimal userland environments such as those generated by Buildroot. Unlike VMware and VirtualBox, QEMU can boot directly from a raw kernel image (e.g., \texttt{bzImage}) without requiring a full operating system or bootloader. This capability is essential for testing custom kernels, isolating kernel-level behavior, and integrating tightly with lightweight Linux distributions.

Second, QEMU enables precise automation and scripting. Every aspect of VM configuration—from CPU model and memory size to disk images, serial interfaces, and network devices—can be fully defined through the command line. This makes QEMU highly suitable for automated, repeatable testing workflows where container configurations and environments must be instantiated deterministically. In contrast, VMware and VirtualBox depend heavily on preconfigured graphical management interfaces and require more overhead to script, which impedes reproducibility and scalability.

Third, QEMU is architecture-agnostic. It supports a wide range of platforms including x86, ARM, MIPS, RISC-V, and PowerPC, enabling cross-architecture security testing. Although Houdini currently focuses on x86\_64, the ability to extend to other architectures in future research is a major benefit. VMware and VirtualBox, by contrast, are restricted to x86/x86\_64 virtualization and do not support full system emulation for other architectures.

Finally, QEMU enhances host safety and test isolation. By running container experiments in a fully virtualized environment, Houdini prevents any malicious or misconfigured container from affecting the host system. Even in the event of a container escape, the damage is confined to the virtual machine, which can be discarded and restarted. This isolation guarantees that all tests are both safe and contained.

For these reasons, QEMU serves as the ideal virtualization backend for Houdini, enabling rigorous, repeatable, and isolated evaluation of container security mechanisms under a wide range of kernel and runtime configurations.

\section{Interface and User Experience Design}

Houdini provides a command-line interface (CLI) as its primary mode of interaction. The CLI is designed to balance flexibility with usability, allowing both simple and advanced users to manage trick execution workflows efficiently. Common commands include \texttt{houdini run} to execute a trick, \texttt{houdini matrix} to expand a configuration file into concrete test cases, and \texttt{houdini rerun} to re-execute failed or inconclusive tests by test ID. Each command provides contextual feedback, including progress indicators, error summaries, and log file locations.

The experiment configuration itself is defined declaratively using a YAML-based format that is human-readable and easy to edit. Users can specify default settings at the global level and override them for individual test cases as needed. This model reduces redundancy and simplifies management of large matrices. 

Test results are presented in a structured, color-coded table that highlights pass/fail status, runtime versions, and key configuration flags. Logs are accessible by test ID and organized in a predictable directory structure, allowing quick navigation and review. Overall, Houdini’s interface is designed to support transparency, traceability, and low-friction experimentation for researchers and engineers.



\section{Runtime Injection via OverlayFS}

Rather than relying on Buildroot’s package infrastructure to download and install Docker and runc inside the VM, we used an overlayfs directory to inject precompiled binaries of both components directly into the root filesystem. This approach provides full control over the runtime environment and ensures that only explicitly specified versions are used during testing. We chose not to use Buildroot packages because doing so introduced limitations: several historical versions of Docker and runc were either unavailable or failed to build consistently within Buildroot’s package ecosystem. By manually sourcing and verifying the binaries, we were able to include a broader range of versions and maintain compatibility across all tests.

\section{Preventing Environmental Drift}

Environmental drift refers to unintentional changes in the system that can affect the outcome of container-based tests. These changes may include automatic updates to the operating system, differences in kernel versions, updated system libraries, or changes in container runtime behavior. Even when a test script remains the same, such background variation can cause container behavior to differ between runs, leading to inconsistent or misleading results.

This is a critical concern for security testing, where consistency is essential. Without environmental stability, researchers may encounter false positives, where a trick appears to succeed because of an incidental change, or false negatives, where a real vulnerability is masked by a system-level patch or configuration shift. In either case, the test results become untrustworthy, and any security claims based on those results may be invalid.

To address this, Houdini is designed to eliminate sources of environmental drift by creating a fully version-controlled, deterministic testing environment. All critical components that influence container behavior, including the Linux kernel, Docker engine, and runc runtime are manually sourced, precompiled, and injected into the root filesystem using an overlay filesystem. This ensures that Houdini uses only explicitly specified versions of each component, independent of host updates or package manager behavior.

This design choice was necessary because many historical versions of Docker and runc are either unavailable through Buildroot’s package system or fail to build due to missing dependencies or upstream changes. By bypassing the package infrastructure and directly injecting verified binaries, Houdini achieves a consistent and reproducible runtime environment across all test sessions.

Additionally, Houdini uses Buildroot to generate a minimal, static guest operating system that excludes background services, update agents, or untracked dependencies that could introduce nondeterminism. The virtual machine is launched without networking to prevent environmental noise from DNS, package mirrors, or remote services. Only a single, tightly scoped communication channel is exposed to the host through a QEMU user-mode forwarded port used for reporting results.

Together, these design decisions ensure that Houdini provides a frozen, deterministic environment where test outcomes are reproducible over time. This stability is especially important for longitudinal analysis, vulnerability tracking, patch verification, and reproducible research, where changes in the results must reflect actual differences in the tested logic, not unintended changes in the surrounding system.

\section{Support for Custom Kernels}

Houdini supports the use of custom Linux kernel builds, enabling researchers to evaluate container behavior on non-standard or experimental kernels. This is particularly useful for testing scenarios that require kernel features not available in mainline releases, or for reproducing issues on distribution-specific kernels.

Users can specify a custom kernel image (bzImage) to be used when launching the QEMU virtual machine. Instead of selecting a precompiled kernel from Houdini’s internal cache, the user provides a path to an externally built kernel through the CLI or configuration file. Houdini will boot the VM with this specified kernel while keeping the rest of the guest environment unchanged.

This feature allows full control over the kernel used during testing, including:

Custom-compiled kernels with modified configuration options.

Vendor-patched or downstream distribution kernels.

Legacy kernels for historical reproduction.

The guest root filesystem, Docker engine, and runc version remain consistent regardless of the selected kernel, ensuring that observed behavior differences are attributable solely to the kernel in use. This modularity supports high-fidelity experimentation without altering the test orchestration pipeline.

To preserve safety, custom kernels are still booted inside isolated QEMU virtual machines. This ensures that even if the custom kernel is unstable or includes regressions, it cannot affect the host system.

This flexibility makes Houdini a valuable platform for kernel-centric container security research, where testing across a range of kernel versions and configurations is essential.

\section{Multi-Runtime Support}

Houdini is designed to support multiple container runtimes, recognizing that the behavior of container isolation features may vary significantly between runtime implementations. While Docker commonly defaults to \texttt{runc}, the Open Container Initiative (OCI) standard permits alternative runtimes that comply with the OCI runtime specification. To ensure comprehensive testing coverage and enable comparisons across runtimes, Houdini includes support for both \texttt{runc} and \texttt{crun}, with a modular design that facilitates the integration of additional runtimes in the future.

Each runtime is precompiled and stored in Houdini’s local binary cache, just like Docker and kernel versions. When executing a test, Houdini injects the selected runtime into the guest virtual machine using its overlay filesystem. The trick configuration specifies the desired runtime, and Houdini ensures that only the corresponding binary is present in the guest environment during test execution. This guarantees that observed behavior reflects the exact runtime under test and is not influenced by any system defaults.

The current implementation includes support for \texttt{runc}, the default Docker runtime, and \texttt{crun}, a high-performance alternative developed by the Open Container Initiative. These runtimes can be swapped per test case, allowing users to evaluate how different runtime implementations handle security boundaries such as capabilities, seccomp profiles, or namespace interactions.

The runtime selection is fully configurable via Houdini’s YAML and Python-based configuration models. In matrix mode, users can easily define a list of runtimes to test, and Houdini will generate test cases across all specified versions. For example, users can observe how a trick behaves when executed under both \texttt{runc} and \texttt{crun}, using the same container image, kernel version, and security flags.

The system is structured to allow the addition of other OCI-compliant runtimes, such as \texttt{youki} or \texttt{Kata Containers}, with minimal changes. New runtimes can be added to the local cache and referenced in trick configurations just like existing ones. This extensibility ensures that Houdini remains adaptable to changes in the container runtime landscape and can be used to validate the correctness and security of emerging runtime implementations.

By supporting multiple runtimes, Houdini empowers researchers to perform comparative analysis, detect behavioral regressions, and evaluate implementation differences that may affect container confinement. This capability is critical for ensuring that container security claims hold consistently across the diverse and evolving set of OCI runtimes.


\section{Limitations}

While Houdini offers a robust foundation for reproducible container security testing, it currently has several limitations. First, it does not support parallel execution; tests are run sequentially to guarantee isolation and avoid concurrency-related nondeterminism. This serialized execution model simplifies debugging and analysis but can be a bottleneck for very large-scale evaluations. Second, Houdini’s runtime support is presently limited to Docker and runc. While the system is extensible, it does not yet include out-of-the-box support for other runtimes like \texttt{crun} or \texttt{youki}, which limits its coverage of the broader container landscape. Finally, Houdini assumes a Linux-based host system with QEMU and KVM support, which may restrict portability to some environments such as macOS or Windows. These limitations represent areas for future development as the framework matures.

\end{doublespace}